\documentclass{article}
\usepackage{pathfinder-readable}

\title{Paid Price Advice in a Double Auction: A Conditional Benchmark}

\begin{document}
\maketitle

\begin{abstract}
This account connects Zhao and colleagues' survey of learning-augmented algorithms with a proposed study
of token costs in an LLM double auction. The peers looked for a way to judge whether paid price advice
improves gains from trade after its inference cost. On the study's fixed value schedule, they found a
conditional, cost-free benchmark that reaches the gross surplus ceiling, and a possible rationing path
where an arbitrarily small price error loses a fixed amount of surplus. These are deductions from the
stated schedule, not measured auction outcomes; the thread ended as a draft because execution rules and
the proposed comparison still need checking.
\end{abstract}

\section{Introduction}

Zhao and colleagues' survey, the supplied paper Q, explains how an algorithm can use a prediction while
still stating what happens when the prediction is wrong \cite{zhao2026learning}. It distinguishes the
thing predicted, the measure of its error, and the information available when the prediction is made. It
also asks what happens when obtaining advice has a cost. Its charged-prediction ratio concerns a
cost-minimising algorithm and an uncharged optimum; it does not directly describe an auction's gains from
trade.

The supplied paper P proposes charging language-model traders for the tokens they generate while deciding
whether and how to quote in a double auction \cite{pathfinder2026price}. Buyers and sellers have private
reservation values, and a crossing quote can execute a trade. The paper uses a fixed, symmetric value
schedule to calculate possible gains from trade. It proposes comparisons with and without a trader-facing
charge, but reports no combined experiment. P calls a different paper, by Struski and colleagues, ``Q''
\cite{struski2026competitive}; that name does not refer to the supplied survey.

The peers used Q's framework to ask what a paid price recommendation would have to beat, what price error
would say about welfare, and what information the adviser could use. Advice was an extension of P, which
specifies costly trader decisions but no separate predictor.

\section{What was done}

The peers first separated Q's cost per predictor call from P's cost for traders' generated tokens. They
used an additive welfare shortfall for P, whose objective differs from Q's cost ratio.

They then checked P's fixed schedule. Its five profitable matches have gains of \(2.5\), \(2\), \(1.5\),
\(1\), and \(0.5\); the next match adds nothing. The peers constructed a scripted price policy with no
language-model calls and checked its outcome if all five pairs traded.

Next they tested a proposed measure of advice quality: the distance between an advised price and the
clearing price. A price just above \(2\) admits an extra seller. Under a possible order of execution,
that seller can take the place of a lower-cost seller, even though the price error is arbitrarily small.
The peers also checked the complementary case: if changing the price changes neither the eligible buyers
nor the eligible sellers, and all of them trade, the price itself does not change gross surplus.

Finally they raised two design objections. A trader's token tariff can change behaviour, while the
analyst's conversion of token use into welfare units measures resource cost. The adviser and a scripted
comparator must also see the same public signals and permitted private information. Neither point was
tested experimentally.

\section{Results}

Let \(G\) be realised gross gains from executed trades, \(C\) the cost of generated tokens under P's
accounting, and \(\rho\geq 0\) the stated conversion of that cost into the same units as gains. Let
\(G^{\star}=7.5\) be the fixed schedule's gross ceiling. The peers' net welfare and additive shortfall
are
\[
W_{\rho}=G-\rho C,
\qquad
R_{\rho}=G^{\star}-W_{\rho}=7.5-G+\rho C.
\]
This is an accounting identity. The symbol \(\rho\) replaces P's conversion symbol because Q uses that
symbol for an unrelated trust setting.

The first result is a conditional benchmark. P's five buyers valued above \(2\) have values totalling
\(13.75\). Its five sellers with costs below \(2\) have costs totalling \(6.25\). A scripted buyer could
bid \(2+\varepsilon\), and a scripted seller could ask \(2-\varepsilon\), when each has a strictly
profitable trade, where \(0<\varepsilon<0.25\); the other traders could abstain. If the auction executes
all five eligible pairs, their order and pairing do not affect the total:
\[
G=13.75-6.25=7.5.
\]
The script uses no generated language-model tokens. On that completed path it attains \(W_{\rho}=7.5\).
No policy on the same values can exceed \(G=7.5\); if \(\rho>0\) and a paid policy uses \(C>0\), its net
welfare is strictly lower. The peers independently checked the sums and construction. P's random
scheduler and finite call cap do not ensure that the script completes the five trades.

The second result concerns price error under rationing. In a posted-price response to P's same values,
price \(2\) makes five buyers and five sellers strictly eligible. At an advised price \(2+\varepsilon\),
where \(0<\varepsilon<0.25\), the seller whose cost is \(2\) also becomes eligible, while the eligible
buyer set stays the same. If that seller trades and the seller whose cost is \(0.75\) does not, then the
seller costs rise from \(6.25\) to \(7.5\). Gross surplus falls to \(6.25\), a loss of \(1.25\), even as
\(\varepsilon\) tends to zero. The peers checked this arithmetic and described a possible crossing-order
path. They did not establish that P's actual auction takes that path.

There is also a simple sufficient condition for avoiding this particular loss. If the advised price
leaves both eligible sets unchanged, both sets have the same number of traders, and every eligible trader
executes once, gross surplus equals the sum of their buyer values less the sum of their seller costs.
Pairing and price then cancel from that sum. This is a statement about full execution with fixed
participation, not a guarantee for P's auction.

\section{What did not hold}

Q's prediction guarantee did not transfer to P. Q's charged ratio is for cost minimisation against an
uncharged positive optimum. P seeks to maximise gains after token costs and does not define a separate
prediction object or prediction error. The additive shortfall above keeps the accounting clear, but
proves no consistency or robustness guarantee.

Small numerical price error did not by itself control welfare. The possible loss of \(1.25\) depends on
which seller trades, so a useful guarantee would need assumptions about allocation, arrival order, or the
margin around the advised price. The example does not show that language-model traders would follow the
proposed quotes.

The tariff question remained open. Write \(\tau\) for the charge a trader faces per token. Changing
\(\tau\) may change both \(G\) and \(C\), while welfare at a fixed resource conversion remains
\(G(\tau)-\rho C(\tau)\). P already proposes declaring a range of welfare conversions, so this
distinction sharpens a future comparison rather than correcting a demonstrated error. The adviser must
also be restricted to information available when advice is issued; giving it the full private value
vector would make a comparison with an ordinary trader unfair. The peers had no run that resolved these
behavioural and information questions.

\section{Why the thread stopped}

The verifier left the thread at \emph{DRAFT} after one round. It accepted the cost-free benchmark and the
small-error rationing example as useful conditional results, and noted that the account kept them
separate from observed behaviour. It also identified execution and stopping rules that still require
implementation checks. The smallest immediate step is to check the released auction's crossing,
repeated-order, and stopping behaviour against the proposed script and rationing path. A decisive next
study would then compare a no-query policy with the same policy offered optional, typed price advice on
prespecified variable value schedules. Both arms would need the same information and scheduler, with
gross surplus, tokens, participation, and executions recorded. P's fixed schedule would serve as a
negative control; no gain from paid advice has yet been shown.

\bibliographystyle{plainurl}
\bibliography{references}

\end{document}
